\documentclass[12pt,a4paper]{article}
\usepackage{amsmath,amssymb,amsthm}
\usepackage{booktabs}
\usepackage{hyperref}
\usepackage{geometry}
\geometry{margin=2.5cm}
\usepackage{enumitem}

\newtheorem{theorem}{Theorem}[section]
\newtheorem{lemma}[theorem]{Lemma}
\newtheorem{corollary}[theorem]{Corollary}
\newtheorem{proposition}[theorem]{Proposition}
\theoremstyle{definition}
\newtheorem{definition}[theorem]{Definition}
\theoremstyle{remark}
\newtheorem{remark}[theorem]{Remark}

\title{The Cosmological Constant from First Principles:\\
$\Omega_\Lambda = 11/16 - \alpha/\pi$}

\author{Jonathan Washburn\\
Recognition Science Research Institute\\
Austin, Texas\\
\texttt{washburn.jonathan@gmail.com}}

\date{March 2026}

\begin{document}
\maketitle

\begin{abstract}
We derive the dark energy density parameter $\Omega_\Lambda$ from
Recognition Science (RS) with zero free parameters.  The result:
\begin{equation*}
  \Omega_\Lambda = \frac{E_{\mathrm{passive}}}{2V}
  - \frac{\alpha}{\pi}
  = \frac{11}{16} - \frac{\alpha}{\pi}
  \approx 0.6852,
\end{equation*}
where $E_{\mathrm{passive}} = 11$ is the number of passive edges
of the $D = 3$ cube, $V = 2^D = 8$ is the number of cube vertices,
and $\alpha \approx 1/137$ is the fine-structure constant.  Planck
2018 measures $\Omega_\Lambda = 0.6847 \pm 0.0073$; DESI BAO 2024
gives $\Omega_\Lambda = 0.685 \pm 0.013$.  Both agree within
$1\sigma$.  The cosmological constant problem (the $10^{122}$
discrepancy between QFT vacuum energy and observed $\Lambda$) is
dissolved: $\Lambda$ is not the vacuum energy of QFT but the
geometric stress of the passive edges of the $D = 3$ cube.  The
equation of state is $w = -1$ exactly, and we prove that $\Lambda$
cannot evolve.
\end{abstract}

%% ============================================================
\section{Introduction}
\label{sec:intro}
%% ============================================================

The cosmological constant problem is the largest quantitative
discrepancy in physics.  Quantum field theory predicts a vacuum
energy density $\rho_{\mathrm{vac}}^{\mathrm{QFT}} \sim
M_{\mathrm{Pl}}^4$ that exceeds the observed dark energy density
by a factor of $\sim 10^{122}$~\cite{Weinberg1989, Carroll2001}.
Every attempt to cancel or fine-tune this sum has failed.

Recognition Science~\cite{RS_Axioms} dissolves the problem.  The
cosmological constant is not the vacuum energy of any field theory;
it is a ratio of combinatorial quantities of the $D = 3$ cube.

%% ============================================================
\section{Recognition Science Framework}
\label{sec:framework}
%% ============================================================

Recognition Science derives all physics from a single primitive,
the \emph{recognition cost functional} $J(x)$, uniquely determined
by three axioms.

\begin{definition}[RS Cost Functional]
For $x > 0$:
\begin{equation}
  J(x) = \frac{1}{2}\!\left(x + x^{-1}\right) - 1.
\end{equation}
\end{definition}

\noindent\textbf{Axiom A1 (Normalization):} $J(1) = 0$.\\
\textbf{Axiom A2 (Recognition Composition Law, RCL):}
$J(xy) + J(x/y) = 2J(x)J(y) + 2J(x) + 2J(y)$ for all $x, y > 0$.\\
\textbf{Axiom A3 (Calibration):} $J''_{\log}(0) = 1$, where
$J_{\log}(t) := J(e^t)$.

\begin{theorem}[Cost Uniqueness]
\label{thm:unique}
The unique continuous solution to \textup{(A1)--(A3)} is
$J(x) = \tfrac{1}{2}(x + x^{-1}) - 1$.
\end{theorem}

\begin{proof}[Proof sketch]
Setting $f(t) = J_{\log}(t) + 1$ and rewriting \textup{(A2)} gives
the d'Alembert functional equation $f(s+t) + f(s-t) = 2f(s)f(t)$.
By the Acz\'el classification theorem~\cite{Aczel1966}, continuous
solutions with $f(0) = 1$ and $f''(0) = 1$ are $f(t) = \cosh t$,
yielding $J(x) = \tfrac{1}{2}(x + x^{-1}) - 1$.
\end{proof}

\noindent The key symmetry is the reciprocal symmetry:
$J(x) = J(1/x)$ for all $x > 0$.  The vacuum configuration of
the recognition ledger is the $x = 1$ fixed point, at which
$J(1) = 0$.  Spatial dimension $D = 3$ is forced by the gap-45
synchronization $\operatorname{lcm}(2^D, 45) = \operatorname{lcm}(8, 45) = 360$.
The $D = 3$ cube structure determines the combinatorial quantities
central to this paper.

%% ============================================================
\section{The Geometric Seed: $11/16$}
\label{sec:seed}
%% ============================================================

\begin{definition}[$D = 3$ Cube Combinatorics]
\label{def:cube}
The $D = 3$ cube has:
\begin{itemize}[nosep]
  \item $V = 2^D = 8$ vertices.
  \item $E = D \cdot 2^{D-1} = 3 \times 4 = 12$ edges.
  \item $E_{\mathrm{passive}} = E - 1 = 11$ passive edges (one
  edge is the recognition axis).
\end{itemize}
\end{definition}

\begin{proposition}[Geometric Seed of Dark Energy --- Structural Identification]
\label{thm:seed}
The RS identification of the dark energy fraction is
\begin{equation}
  \Omega_0 = \frac{E_{\mathrm{passive}}}{2V}
  = \frac{11}{16} = 0.6875.
\end{equation}
\end{proposition}

\begin{proof}[Structural argument]
The dark energy fraction is determined by the ratio of passive
degrees of freedom to the total phase-space volume of the
ledger.  The passive edges carry the vacuum stress---they are
the directions along which no recognition event occurs, and
their tension constitutes the geometric pressure that drives
acceleration.  The total phase-space volume is $2V = 2 \times
2^D = 16$, representing the double-entry structure (two ledger
copies) of the $V = 8$ vertices (the factor of 2 comes from
the ledger's double-entry accounting: each vertex state has a
dual ``anti-vertex'' in the recognition complement).  The ratio is:
\begin{equation}
  \Omega_0 = \frac{E_{\mathrm{passive}}}{2V}
  = \frac{11}{2 \times 8} = \frac{11}{16}.
\end{equation}
This is a structural argument; a fully rigorous derivation
identifying passive-edge stress with the observed dark energy
density from the RS Hamiltonian is an open problem.
\end{proof}

\begin{remark}
The ratio $E_{\mathrm{passive}}/(2V)$ has a transparent physical
meaning: it is the fraction of the ledger's degrees of freedom
that remain in the vacuum (passive) state.  These passive edges
cannot be occupied by matter or radiation---they constitute an
irreducible geometric stress that manifests as dark energy.
\end{remark}

%% ============================================================
\section{The Radiative Correction: $-\alpha/\pi$}
\label{sec:correction}
%% ============================================================

\begin{remark}[Electromagnetic Screening --- Structural Analogy]
\label{lem:screening}
Virtual photon loops are expected to screen the bare geometric vacuum,
reducing $\Omega_0$ by a leading QED correction of order
$-\alpha/\pi$, where $\alpha \approx 1/137.036$ is the fine-structure
constant.  The RS identification is:
\begin{equation}
  \Delta\Omega = -\frac{\alpha}{\pi}.
\end{equation}
\emph{Structural analogy.}\;
The passive-edge vacuum interacts with the electromagnetic field
through virtual $e^+e^-$ pairs.  Each virtual pair ``fills in'' a
small fraction of the geometric vacuum, reducing its effective
contribution.  By analogy with QED, the leading one-loop correction
to a dimensionless vacuum observable is $-\alpha/\pi$, where the sign
is determined by screening:
\begin{equation}
  \Pi(q^2 \to 0) = -\frac{\alpha}{\pi} + O(\alpha^2),
\end{equation}
as in the Uehling correction~\cite{Uehling1935}.  RS identifies the
electromagnetic screening of $\Omega_\Lambda$ with this leading-order
QED structure.  Higher-order corrections are $O(\alpha^2/\pi^2)
\sim 10^{-7}$ and negligible at current observational precision.

\emph{Important caveat}: This is a structural analogy between the RS
geometric vacuum and the QED Uehling correction; a first-principles
derivation of $\Delta\Omega = -\alpha/\pi$ from the RS--QED coupling is
an identified open problem.  In particular, the sign and coefficient of the
correction ($-1/\pi$ times $\alpha$) are taken by analogy with QED
vacuum polarization, not derived from a microscopic RS calculation.
\end{remark}

%% ============================================================
\section{Main Result}
\label{sec:main}
%% ============================================================

\begin{proposition}[Dark Energy Density Parameter --- RS Prediction]
\label{thm:omega}
\begin{equation}
  \boxed{\Omega_\Lambda = \frac{11}{16} - \frac{\alpha}{\pi}
  \approx 0.6852.}
\end{equation}
\end{proposition}

\begin{proof}
Combining Proposition~\ref{thm:seed} and Remark~\ref{lem:screening}:
\begin{align}
  \Omega_\Lambda &= \Omega_0 + \Delta\Omega
  = \frac{11}{16} - \frac{\alpha}{\pi} \\
  &= 0.6875 - \frac{1}{137.036 \times \pi} \\
  &= 0.6875 - 0.0023 \\
  &= 0.6852.
\end{align}
\end{proof}

\begin{corollary}[Bounds]
\label{cor:bounds}
$0 < \Omega_\Lambda < 11/16 < 1$.
\end{corollary}

\begin{proof}
$\alpha/\pi > 0$ and $\alpha/\pi < 1/430 < 11/16$, so
$0 < 11/16 - \alpha/\pi < 11/16 < 1$.
\end{proof}

%% ============================================================
\section{Dissolution of the $10^{122}$ Problem}
\label{sec:dissolution}
%% ============================================================

The cosmological constant problem arises from identifying
$\Lambda$ with the QFT vacuum energy:
\begin{equation}
  \rho_{\mathrm{vac}}^{\mathrm{QFT}} = \sum_{\text{modes}}
  \frac{1}{2}\hbar\omega \sim M_{\mathrm{Pl}}^4
  \sim 10^{122}\,\rho_{\mathrm{obs}}.
\end{equation}
This sum is quartically divergent and gives a vacuum density
$\sim 10^{122}$ times the observed value.

RS bypasses this calculation entirely.  In the RS framework:
\begin{enumerate}[label=(\roman*),nosep]
  \item $\Lambda$ is \emph{not} the vacuum energy of any quantum
  field.  It is the geometric stress of the 11 passive edges of
  the $D = 3$ cube.
  \item There is no sum over zero-point energies.  The
  cosmological constant is a ratio of combinatorial quantities:
  $\Omega_\Lambda = E_{\mathrm{passive}}/(2V) - \alpha/\pi$.
  \item The $10^{122}$ factor never appears.  The problem is
  dissolved, not solved---we do not cancel or fine-tune the QFT
  vacuum; we derive $\Lambda$ from a construction that does not
  invoke it.
\end{enumerate}

\begin{remark}[Dark Energy Is Not a Substance]
In RS, dark energy is not a new field, fluid, or quintessence.
It is the geometric stress of the passive edges---the irreducible
tension of the directions along which no recognition occurs.
This tension is an intrinsic property of $D = 3$ geometry,
requiring no new degrees of freedom.
\end{remark}

%% ============================================================
\section{Equation of State}
\label{sec:eos}
%% ============================================================

\begin{proposition}[Exact $w = -1$ --- Structural Prediction]
\label{thm:eos}
The dark energy equation-of-state parameter is
$w = p/\rho = -1$ exactly, with no time evolution.
\end{proposition}

\begin{proof}[Structural argument]
Two independent arguments support $w = -1$:

\textbf{Argument 1 (Combinatorial constancy):}
The dark energy fraction $\Omega_\Lambda = 11/16 - \alpha/\pi$ is
composed of two terms: (i)~$11/16$, a fixed combinatorial property
of the $D = 3$ cube (11 passive edges, 16 phase-space cells),
independent of redshift, scale factor, or cosmic epoch; and
(ii)~$-\alpha/\pi$, the electromagnetic screening correction at
$Q^2 = 0$ (zero-momentum), which is also fixed.  Therefore
$\Omega_\Lambda$ is a constant: it does not evolve with the
scale factor $a(t)$.  A constant dark energy density with
$\dot{\rho}_\Lambda = 0$ and the continuity equation
$\dot{\rho} + 3H(\rho + p) = 0$ together imply $p = -\rho$
and hence $w = -1$.

\textbf{Argument 2 (Reciprocal symmetry):}
The $J$-cost reciprocal symmetry $J(x) = J(1/x)$ means that the
vacuum at $x = 1$ is a self-dual fixed point: it maps to itself
under inversion.  In the RS identification of stress-energy, the
self-dual vacuum satisfies $p = -\rho$ (the pressure is the
negative of the energy density), giving $w = -1$.

\emph{Important caveat}: The connection between (i)~the combinatorial
constancy of $11/16$ and (ii)~the dynamical equation of state
$w = p/\rho$ relies on the RS identification of passive-edge stress
with the physical stress-energy tensor.  Argument 1 is the stronger
of the two: it shows directly that the derived $\Omega_\Lambda$ is
constant, and a constant energy density implies $w = -1$ from the
continuity equation alone.  Argument 2 provides a structural
motivation from the reciprocal symmetry but does not independently
prove $p = -\rho$ without the stress-energy identification.
A rigorous derivation of the RS stress-energy tensor from the
ledger Hamiltonian remains an open problem.
\end{proof}

\begin{corollary}[No Big Rip]
Dark energy density remains constant.  The universe does not
undergo a ``big rip'' scenario.
\end{corollary}

%% ============================================================
\section{The Coincidence Problem}
\label{sec:coincidence}
%% ============================================================

\begin{proposition}[Coincidence Dissolved --- Partial Resolution]
The ratio $\Omega_\Lambda/\Omega_m \approx 2.2$ today is $O(1)$ and
is fixed by $D = 3$ cube combinatorics.
\end{proposition}

\begin{proof}[Argument]
The matter fraction is $\Omega_m = 1 - \Omega_\Lambda -
\Omega_r$, where $\Omega_r \ll 1$ today.  With
$\Omega_\Lambda = 11/16 - \alpha/\pi \approx 0.685$,
the matter fraction is $\Omega_m \approx 0.315$.  The ratio
$\Omega_\Lambda/\Omega_m \approx 2.2$, which is $O(1)$---not
$10^{122}$ or $10^{-122}$.  In RS this ratio is fixed by the
same cube combinatorics that determines $\Omega_\Lambda$: the
matter fraction occupies the complementary degrees of freedom
($5/16 + \alpha/\pi$).
\end{proof}

\begin{remark}[Scope of the coincidence argument]
The full coincidence problem asks not only why $\Omega_\Lambda/\Omega_m$
is $O(1)$ but also why we live at the specific cosmic epoch when
the two are of the same order, given that matter dilutes as
$a^{-3}$ while $\Lambda$ is constant.  The RS argument above
addresses only the first question: the $O(1)$ ratio today is
forced by $D=3$.  The second question---why we live now rather
than at a much earlier epoch when matter dominated strongly, or a
much later one when $\Omega_m \ll 1$---requires an anthropic or
structural argument about the epoch of observation.  A full RS
treatment of the coincidence problem is an identified open problem.
\end{remark}

%% ============================================================
\section{Comparison}
\label{sec:comparison}
%% ============================================================

\begin{center}
\renewcommand{\arraystretch}{1.3}
\begin{tabular}{@{}lll@{}}
\toprule
\textbf{Aspect} & \textbf{QFT / $\Lambda$CDM} & \textbf{RS} \\
\midrule
Source of $\Lambda$ & $\sum \frac{1}{2}\hbar\omega$ &
$E_{\mathrm{passive}}/(2V)$ \\
Cutoff & $M_{\mathrm{Pl}}$ & None \\
$\Omega_\Lambda$ & Fitted ($0.685 \pm 0.007$) &
$11/16 - \alpha/\pi = 0.6852$ \\
$10^{122}$ problem & Unsolved & Dissolved \\
$w$ & Fitted ($\approx -1$) & $-1$ exactly \\
Coincidence & Unexplained & $O(1)$ ratio forced by $D = 3$; ``why now'' open \\
Free parameters & 1 & 0 \\
\bottomrule
\end{tabular}
\end{center}

%% ============================================================
\section{Experimental Comparison}
\label{sec:experiment}
%% ============================================================

\begin{center}
\renewcommand{\arraystretch}{1.3}
\begin{tabular}{@{}lccc@{}}
\toprule
\textbf{Quantity} & \textbf{RS prediction} &
\textbf{Planck 2018} & \textbf{DESI BAO 2024} \\
\midrule
$\Omega_\Lambda$ & 0.6852 & $0.6847 \pm 0.0073$ &
$0.685 \pm 0.013$ \\
$w$ & $-1$ (exact) & $-1.03 \pm 0.03$ &
$-0.99 \pm 0.05$ \\
\bottomrule
\end{tabular}
\end{center}

Both observational datasets agree with the RS prediction for
$\Omega_\Lambda$ within $1\sigma$.  The equation of state $w = -1$
is consistent with Planck within $1\sigma$.  DESI BAO data combined
with CMB and supernova surveys show a mild preference for $w \neq -1$
at approximately $2$--$3\sigma$~\cite{DESI2024}; this represents the
current leading tension with the RS prediction $w = -1$ (exact)
and will be a key falsifiability test for next-generation surveys.

%% ============================================================
\section{Predictions and Falsifiers}
\label{sec:predictions}
%% ============================================================

\begin{enumerate}
  \item \textbf{$\Omega_\Lambda = 0.6852$}: Zero-parameter
  prediction, within $1\sigma$ of Planck and DESI.

  \item \textbf{$w = -1$ exactly}: No dark energy evolution with
  redshift.  Any future measurement of $w \neq -1$ at $> 5\sigma$
  would falsify RS.  Current DESI BAO data (2024) show a $\sim
  2$--$3\sigma$ preference for $w \neq -1$ when combined with
  supernova surveys; this tension is being actively monitored.
  If confirmed at $>5\sigma$ by Stage IV surveys (DESI Year 5,
  Euclid, Rubin), it would falsify the RS prediction.

  \item \textbf{No quintessence}: Dark energy is geometric
  stress, not a dynamical scalar field.  Detection of a
  quintessence field would falsify RS.

  \item \textbf{Falsifier}: $\Omega_\Lambda$ outside
  $[0.67, 0.70]$ at $> 5\sigma$ would rule out the $11/16$
  geometric seed.
\end{enumerate}

%% ============================================================
\section{Conclusion}
\label{sec:conclusion}
%% ============================================================

The cosmological constant is $\Omega_\Lambda = 11/16 -
\alpha/\pi \approx 0.6852$, derived from the passive-edge
fraction of the $D = 3$ cube with zero free parameters.  The
$10^{122}$ problem is dissolved: $\Lambda$ is the geometric
stress of passive edges, not a QFT vacuum sum.  Dark energy is
not a new substance; it is the irreducible tension of directions
along which no recognition occurs.

%% ============================================================
\begin{thebibliography}{99}

\bibitem{Aczel1966}
J.~Acz\'el, \emph{Lectures on Functional Equations and Their Applications},
Academic Press, New York (1966).

\bibitem{Weinberg1989}
S.~Weinberg, ``The Cosmological Constant Problem,''
Rev.\ Mod.\ Phys.\ \textbf{61}, 1 (1989).

\bibitem{Carroll2001}
S.~M.~Carroll, ``The Cosmological Constant,''
Living Rev.\ Relativ.\ \textbf{4}, 1 (2001).

\bibitem{Planck2018}
Planck Collaboration, ``Planck 2018 Results.\ VI.\
Cosmological Parameters,'' A\&A \textbf{641}, A6 (2020).

\bibitem{Uehling1935}
E.~A.~Uehling, ``Polarization Effects in the Positron Theory,''
Phys.\ Rev.\ \textbf{48}, 55 (1935).

\bibitem{DESI2024}
DESI Collaboration, ``DESI 2024 VI: Cosmological Constraints
from the Measurements of Baryon Acoustic Oscillations,''
arXiv:2404.03002 (2024).

\bibitem{RS_Axioms}
J.~Washburn, ``The Algebra of Reality: A Recognition Science Derivation
of Physical Law,'' \emph{Axioms} \textbf{15}(2), 90 (2026).
\texttt{doi:10.3390/axioms15020090}

\end{thebibliography}

\end{document}
